\subsection{An algebraic certificate that becomes a quantum experiment}
\label{geo:hesse-cert}

The cubic classification identifies an ensemble from exact equations.
The next result is quantitative and operational. A request of size \(k\)
means that a hidden ray \(z\sim\mu\) is drawn once and \(k\) copies of that
same state are prepared; independent requests use independent draws.
This is a matched-replica model with trusted dimension, basis, and
measurements. Its average request state is \(M_k(\mu)\). Redrawing the
hidden ray separately for every copy would instead reveal only \(M_1\).

Use pure-state trace distance
\[
 d(z,w)=\frac12\|P_z-P_w\|_1
       =\sqrt{1-|\langle z,w\rangle|^2}
\]
and its \(1\)-Wasserstein distance \(W_1\) on ensembles.
In \(\operatorname{Sym}^3\mathbb C^3\), define orthonormal vectors
\[
 \phi_F=\frac{|000\rangle+|111\rangle+|222\rangle}{\sqrt3},\qquad
 \phi_T=\frac1{\sqrt6}\sum_{\pi\in S_3}|\pi(012)\rangle,\qquad
 Q=|\phi_F\rangle\langle\phi_F|+|\phi_T\rangle\langle\phi_T|.
\]
Its coherent-state symbol is
\begin{equation}
 q(z)=\langle z^{\otimes3},Qz^{\otimes3}\rangle
 =\frac13|z_0^3+z_1^3+z_2^3|^2+6|z_0z_1z_2|^2.
 \label{geo:q}
\end{equation}
The two cubic equations cut out precisely the nine rays \(\mathcal H\).

\fig{fig02_cubic_witness}{A positive cubic witness locates the Hesse
ensemble. (a) The exact symbol $q$ on the slice
$z=(0,\cos s,e^{i\theta}\sin s)$ has three isolated zeros at equal
coordinate moduli and the Hesse phases. (b) For 2,500 independent complex
Gaussian rays normalized to unit length (fixed reproducible seed), the
distance to the nearest Hesse ray is plotted against $q$. The solid curve
is the proved bound $\min\{1,\sqrt{27q/4}\}$; the samples illustrate the
bound and do not establish it. The inset shows the same symbol near one
zero as a surface, with the plane $q=0.03$ for scale.}{fig:witness}

\begin{theorem}[Robust Hesse ensemble certificate]
\label{geo:hesse-bound}
For any Borel probability measure \(\mu\) on \(\mathbb P^2\), put
\[
 \epsilon_Q=\operatorname{tr}(QM_3(\mu)),\qquad
 \delta_1=\|M_1(\mu)-I_3/3\|_{\rm HS},\qquad K=1+\frac3{\sqrt2}.
\]
Then
\begin{equation}
 W_1(\mu,\mu_{\mathcal H})
 \le K\sqrt{\frac{27}{4}\epsilon_Q}+\frac32\delta_1.
 \label{geo:robust}
\end{equation}
In particular, \(M_1(\mu)=I_3/3\) and \(\epsilon_Q=0\) characterize
\(\mu_{\mathcal H}\) among all probability ensembles.
\end{theorem}

\begin{proof}
First we prove the global error bound
\begin{equation}
 d(z,\mathcal H)^2\le\frac{27}{4}q(z).
 \label{geo:pointwise}
\end{equation}
For a normalized binary vector \((u,v)\), let \(\Delta\le\pi/3\) be its
phase error from the nearest vector \((1,-\omega^j)/\sqrt2\).
Put \(t=2|uv|\), \(y=t\cos\Delta\). Its squared distance to this
three-point set is \((1-y)/2\), whereas
\[
 |u^3+v^3|^2
 =(1-y)\left(1+y+y^2-\frac34t^2\right).
\]
The second factor is at least one, since
\[
 y+y^2-\frac34t^2
 =(y-t/2)(y+t/2+1)+\frac12t(1-t)\ge0.
\]
Thus binary squared distance is at most \(|u^3+v^3|^2/2\).

For a qutrit let \(a\ge b\ge c\) be its squared coordinate moduli and
\(f=|\sum_i z_i^3|\). Delete the minimum coordinate and use the binary
estimate. Since \(c\le1/3\),
\[
 d(z,\mathcal H)^2
 \le c+\frac{(f+c^{3/2})^2}{2(1-c)^2}
 \le\frac54c+\frac94f^2.
\]
If \(a\le3/4\), then \(ab\ge3/32\), and this is at most
\(40abc/3+9f^2/4\le(27/4)q(z)\).
If \(a\ge3/4\), then
\[
 f\ge a^{3/2}-(1-a)^{3/2}\ge2a-1,\qquad
 d(z,\mathcal H)^2\le\frac{3-a}{4}.
\]
The identity
\(9(2a-1)^2-(3-a)=(4a-3)(9a-2)\ge0\) again proves
\eqref{geo:pointwise}.

Map each ray to a nearest point of \(\mathcal H\), obtaining weights
\(w_i\), and put \(L^2=\int d(z,\mathcal H)^2\,d\mu\).
This transport costs at most \(L\), and the first-moment error of the
transported measure is at most \(\delta_1+\sqrt2L\).
For \(u_i=w_i-1/9\), the SIC Gram matrix gives
\[
 \left\|\sum_i u_iP_i\right\|_{\rm HS}^2=\frac34\sum_i u_i^2.
\]
The weights' total-variation distance from uniform is at most
\(\sqrt3\,\delta_1+\sqrt6\,L\).
All pairwise Hesse distances are \(\sqrt3/2\), so correcting the weights
costs at most \((3/2)\delta_1+(3/\sqrt2)L\).
Combine the two transports and \(L^2\le27\epsilon_Q/4\).
\end{proof}

All measurements in the theorem are positive operators.
The Hesse third moment has the particularly simple form
\begin{equation}
 M_3(\mu_{\mathcal H})=\frac18(I_{10}-Q).
 \label{geo:hesse-third}
\end{equation}
The vectors \(h^{\otimes3}\) form a nine-vector equiangular tight frame
in the eight-dimensional kernel of \(Q\). Algebraically, the ideal
\((z_0^3+z_1^3+z_2^3,z_0z_1z_2)\) is a complete intersection, with
Hilbert series
\[
 \frac{(1-t^3)^2}{(1-t)^3}.
\]
Its Hilbert function starts \(1,3,6,8,9,9,\ldots\), giving the ranks
available to coherent tensors at successive replica numbers.
At three replicas the density matrix has rank eight, yet its
decomposition into pure coherent tensors is uniquely the nine-point
Hesse mixture. Indeed, the coherent mixtures annihilated by \(Q\) are
supported on \(\mathcal H\), and the Gram matrix of their nine
projectors is
\[
 \frac{63I_9+\mathbf1\mathbf1^{\mathsf T}}{64},
\]
so those projectors are linearly independent. They span an exposed
simplex face of the convex set of bosonic separable states.

Three matched replicas are also necessary for distribution-free
identification. Haar measure and \(\mu_{\mathcal H}\) have equal first
and second moments, so any protocol with independent request blocks of
size at most two receives identical quantum states in the two cases,
including joint processing of all blocks. At order three,
Haar has \(M_3=I_{10}/10\) and \(Q\)-probability \(1/5\), whereas Hesse
has \(Q\)-probability zero. The degree of the defining equations has
therefore acquired a direct experimental meaning.

\paragraph{Finite-sample certification.}
On \(N_3\) triple requests measure \(\{Q,I-Q\}\), observing \(s\) clicks.
Let \(q_U\) be the one-sided Clopper--Pearson upper bound at failure
probability \(\alpha/2\). For \(s<N_3\), it is the
\(1-\alpha/2\) quantile of \(\operatorname{Beta}(s+1,N_3-s)\); if
\(s=N_3\), put \(q_U=1\). In particular, zero clicks give
\[
 q_U=1-(\alpha/2)^{1/N_3}.
\]
On \(N_1\) single requests measure the SIC POVM \(\{P_h/3\}\),
with empirical frequencies \(\widehat p_h\), and set
\[
 \widehat\delta_1=\sqrt{12}
 \left\|\widehat p-\frac19\mathbf1\right\|_2,\qquad
 u_{N_1}=
 \frac{\sqrt{12}\bigl(1+\sqrt{\log(2/\alpha)}\bigr)}{\sqrt{N_1}}.
\]
Multinomial mean-square error and bounded differences imply, with
confidence at least \(1-\alpha\),
\begin{equation}
 W_1(\mu,\mu_{\mathcal H})
 \le\min\left\{1,\,
 K\sqrt{\frac{27}{4}q_U}
 +\frac32(\widehat\delta_1+u_{N_1})\right\}.
 \label{geo:confidence}
\end{equation}
The inverse SIC relation is \(M_1=4\sum_h p_hP_h-I\).
For example, the sufficient budgets
\[
 N_3\ge27K^2\varepsilon^{-2}\log(2/\alpha),\qquad
 N_1\ge432\varepsilon^{-2}
                 \bigl(1+\sqrt{\log(2/\alpha)}\bigr)^2
\]
make the acceptance rule \(s=0\),
\(\widehat\delta_1\le\varepsilon/6\) sound and complete for the ideal
Hesse ensemble at confidence \(1-\alpha\). The number of physical
copies is \(3N_3+N_1\).

The accuracy exponent \(\varepsilon^{-2}\) is optimal for fixed request
patterns. For \(0\le\theta\le\pi/6\), rotate the ensemble by
\(U_\theta=\operatorname{diag}(e^{i\theta},e^{-i\theta},1)\). Then
\[
 W_1(\mu_{\mathcal H},U_{\theta*}\mu_{\mathcal H})
 =\frac{\sin\theta+2\sin(\theta/2)}3,\qquad
 \epsilon_Q=
 \frac{\sin^2(3\theta)+2\sin^2(3\theta/2)}{18}.
\]
Coupling hidden labels and using tensor-product fidelity bounds the
trace distance of the two \(n\)-copy experiments by \(\sqrt n\sin\theta\).
Constant-success testing at ensemble distance \(\varepsilon\) thus
requires \(n=\Omega(\varepsilon^{-2})\). This assertion concerns
fixed request patterns; it makes no claim about every adaptive
stopping convention.

\paragraph{A portable principle.}
Let a finite ray set \(S=\{h_1,\ldots,h_N\}\) be cut out by degree-\(k\)
equations whose coefficient space has orthogonal projector \(Q\).
Suppose a certified global bound is available:
\[
 d(z,S)^2\le
 C\langle z^{\otimes k},Qz^{\otimes k}\rangle.
\]
Suppose also that the weight map
\(u\mapsto\sum_j u_jP_{h_j}^{\otimes r}\), restricted to
\(\sum_j u_j=0\), has least singular value \(s_r>0\). If the observed
\(r\)-th moment differs from that of target weights by \(\delta\)
in Hilbert--Schmidt norm, and the \(Q\)-probability is \(\epsilon\),
the same transport argument gives
\begin{equation}
 W_1(\mu,\mu_*)
 \le\sqrt{C\epsilon}
 +\frac{\operatorname{diam}(S)\sqrt N}{2s_r}
       \left(\delta+\sqrt{2rC\epsilon}\right).
 \label{geo:general-cert}
\end{equation}
Equations localize the support; moments condition its weights.
Transverse isolated zeros give local quadratic error bounds. An
effective global \(C\) additionally needs a Jacobian estimate near the
zeros and a positive minimum away from them. These are concrete
targets for exact algebra and certified numerics.
